\chapter{Complete Pick-and-Place Algorithm}
\label{ch:task}

\section{Task target definitions}
The launch arguments are \code{headless} (default false), \code{run_task}
(default true), \code{pickup_mode} (default \code{far}), and \code{seed}
(default 7). The launch code sets the sampled/optimized Box B coordinates into
the task node's \code{pickup_x} and \code{pickup_y} parameters. The node
sets $z=0.025$ m for its pickup center. It sets destination to
$(place_x,place_y,0.225)$ m with defaults $(0.55,0.35,0.225)$.

The goal orientation is fixed to $R_{down}$. The object starts upright check
uses quaternion rotation magnitude and rejects initial deviation above 0.01
rad. The pickup must match the configured pose within \SI{0.01}{m} before
planning. This is a scene-level consistency check; it is not a vision-based
object detection algorithm.

\section{Chronological execution}
\Cref{fig:workflow} shows the decisions and sequencing. Below each phase is
connected to actual class/function names and observable behavior.

\begin{figure}[htbp]
\centering\resizebox{.75\textwidth}{!}{\begin{tikzpicture}[
  node distance=7mm,
  step/.style={draw,rounded corners,fill=blue!5,align=center,text width=49mm,minimum height=9mm},
  decision/.style={draw,diamond,aspect=2,fill=yellow!12,align=center,text width=33mm,inner sep=1pt},
  arr/.style={-{Stealth[length=2mm]},thick}
]
\node[step] (start) {Launch scene; spawn URDF; activate three controllers};
\node[step,below=of start] (feedback) {Wait for clock, arm joints, Box B pose, controller action server};
\node[step,below=of feedback] (detach) {Confirm initial grasp detached; validate pickup pose and upright cube};
\node[step,below=of detach] (preflight) {Compute approach IK; precompute every Cartesian segment};
\node[decision,below=of preflight] (ok) {All IK and path checks pass?};
\node[step,right=19mm of ok] (abort) {Abort before motion; log error};
\node[step,below=of ok] (approach) {Open; joint-space approach; Cartesian descent; check TCP-to-cube error};
\node[step,below=of approach] (grasp) {Close to $-0.162$ rad; attach fixed-joint plugin; confirm state};
\node[step,below=of grasp] (carry) {Lift, inward raise if far, transfer, lower 1 mm above target};
\node[step,below=of carry] (release) {Check placement feedback; detach; open; retreat; settle};
\node[decision,below=of release] (verify) {Fresh final pose within 8 mm and 1 degree?};
\node[step,below=of verify] (done) {Emit \code{PICK AND PLACE VERIFIED}; optional JSON result};
\draw[arr] (start)--(feedback); \draw[arr] (feedback)--(detach); \draw[arr] (detach)--(preflight); \draw[arr] (preflight)--(ok);
\draw[arr] (ok)--node[right,font=\scriptsize]{yes}(approach);
\draw[arr] (ok)--node[above,font=\scriptsize]{no}(abort);
\draw[arr] (approach)--(grasp); \draw[arr] (grasp)--(carry); \draw[arr] (carry)--(release); \draw[arr] (release)--(verify); \draw[arr] (verify)--node[right,font=\scriptsize]{yes}(done);
\draw[arr] (verify.east)-|node[pos=.25,right,font=\scriptsize]{no; fail}(abort.south);
\end{tikzpicture}
}
\caption{Feedback checks and preflight ordering in the task node. All path
segments are solved before the first arm trajectory is sent.}
\label{fig:workflow}
\end{figure}

\begin{enumerate}
 \item \textbf{Start world and robot.} \code{pick_place.launch.py} chooses
 pickup, writes temporary SDF, applies \code{task_description}, starts
 Gazebo/clock bridge/state publisher/spawn, and spawns three controllers. A
 failed controller-spawner process triggers launch shutdown. The task node is
 conditionally started only after the controller process exits successfully.

 \item \textbf{Wait for feedback.} \code{PickPlace.run} waits up to 60 s for
 complete arm joint feedback, the Box B model pose, and a nonzero simulated
 clock. It then waits up to 30 s for the arm trajectory action server. Joint
 feedback is reordered by the exact \code{ARM_JOINTS} names.

 \item \textbf{Reset/validate grasp state.} The node asks the Gazebo plugin to
 detach and confirms the state before motion. It waits for settling, compares
 object position against \code{self.pickup} within 1 cm, stores initial
 quaternion, and checks upright within 0.01 rad.

 \item \textbf{Choose above-pickup pose and waypoints.}
 \code{carry_waypoints(pickup,destination)} decides whether pickup radius
 exceeds 0.80 m. Near pickup goes to z=0.40 m above the cube. Far pickup uses
 z=0.10 m initially and adds the inward waypoint $(0.45,-0.45,0.40)$ m.
 Transfer is at z=0.40 m; lowering ends at destination plus 1 mm in z.

 \item \textbf{Solve and preflight.} The node solves the above-pickup IK with
 deterministic restarts, then generates all waypoint Cartesian segments in
 sequence, seeding each from the preceding endpoint. Any unreachable pose,
 bad starting orientation, or discontinuous joint branch raises before arm
 movement. This preflight order is a useful defense against mid-task IK failure.

 \item \textbf{Open and approach.} \code{fingers(0.15)} confirms a command
 subscriber and live driver feedback, interpolates the opening command, and
 checks the measured result. \code{joint_segment} moves from measured current
 $q$ to the above-pickup configuration. The first Cartesian segment descends
 to Box B.

 \item \textbf{Check grasp point, close, attach.} The current FK grasp-frame
 point must be within \SI{4}{mm} of Box B's center. If not, grasp is refused.
 The node closes the driver to $-0.162$ rad, asks the fixed-joint plugin to
 attach, and waits up to 8 s for \code{attached}. It sets its carrying flag
 only after confirmation.

 \item \textbf{Lift and transfer.} The node executes the remaining path
 segments except the final retreat. These are vertical lift; optional far
 inward/raise; upright transfer; and lowering to z=0.226 m. After each arm
 goal it checks controller success and joint position feedback. While
 carrying, it also checks fresh Box B pose, attached state, and maximum
 sampled orientation error at most 1 degree.

 \item \textbf{Release and retreat.} Before detach, the current Box B center
 must be within 8 mm of destination. The task confirms detach, clears its
 carrying flag, opens the gripper, and executes the last waypoint back to the
 overhead transfer pose. It then waits 2 s in simulated time for settling.

 \item \textbf{Final acceptance.} The latest pose must be fresher than 1 s;
 Euclidean position error must not exceed 8 mm; orientation change from the
 original object rotation must not exceed 1 degree. On success it logs
 \code{PICK AND PLACE VERIFIED}. An optional nonempty \code{result_file}
 parameter writes JSON, but the launch does not set this parameter by default.
\end{enumerate}

\section{Task pseudocode}
\begin{algorithm}[htbp]
\caption{Source-faithful outline of \code{PickPlace.run}}
\label{alg:pick-place}
\begin{algorithmic}[1]
\State wait for joint, Box B pose, simulation clock, and arm action server
\State request and confirm initial detach; validate configured pickup and upright pose
\State $above,waypoints\gets\texttt{carry\_waypoints}(pickup,destination)$
\State $q_{above}\gets\texttt{IK}(above,q_{measured},restart=\texttt{true})$
\State precompute every Cartesian segment, warm-starting each from prior endpoint
\State open gripper; execute quintic joint-space approach
\State execute first segment (descend); reject if grasp frame is farther than 4 mm
\State close driver to $-0.162$ rad; attach object and confirm Gazebo state
\For{each path after descent except final retreat}
  \State send \code{FollowJointTrajectory} action and verify result/feedback
  \State if carrying, check fresh pose, attached state, and sampled orientation error
\EndFor
\State reject if pre-release object center is more than 8 mm from destination
\State detach; open; execute retreat; settle for 2 s
\If{final pose fresh, position error $\leq 8$ mm, angle $\leq1^\circ$}
  \State report verified success and optional JSON
\Else
  \State report failure
\EndIf
\end{algorithmic}
\end{algorithm}

\section{Move helper and ROS messages}
\code{PickPlace.move} takes $(t,qs,vs)$ and a label. It builds a
\code{FollowJointTrajectory.Goal}, assigns six joint names, and for each sample
converts floating-point seconds to integer nanoseconds in
\code{Duration(sec,nanosec)}. It waits for server acceptance and result,
cancels the goal if waiting raises, and rejects non-success action codes. Once
successful, it waits 0.2 s, then verifies measured joints against the final
sample within a max-coordinate error of 0.015 rad. The carrying feedback test
is applied after each move.

\section{Failure behavior and safety scope}
The \code{/gp7/stop} callback sets a cancellation flag and cancels an active
goal. A wait notices the flag and raises, retaining an attached object if the
task was carrying. The outer handler attempts cancellation and exits with
failure. This is a simulation stop/cancel interface; it is not a certified
emergency stop, safe torque off, or physical risk-reduction function. After a
cancellation while attached, the documented safe operational response for
this simulation is to restart the scene before another task.
